% Ficha de estudio 8 del Taller virtual ATA EPE: guardar, compartir, retos y realidad aumentada.
\documentclass[11pt,a4paper]{article}
% Preámbulo común de las fichas de estudio del Taller virtual ATA EPE.
\usepackage[spanish,es-noshorthands,es-tabla]{babel}
\usepackage[utf8]{inputenc}
\usepackage[T1]{fontenc}
\usepackage{lmodern}
\usepackage[margin=2.2cm]{geometry}
\usepackage[table]{xcolor}
\usepackage{amsmath,amssymb,bm}
\usepackage{tabularx,array,booktabs}
\usepackage{enumitem}
\usepackage{fontawesome5}
\usepackage[most]{tcolorbox}
\usepackage[hidelinks]{hyperref}
\usepackage{listings}

\definecolor{morado}{HTML}{2A1B7E}
\definecolor{amarillo}{HTML}{FFCF2E}
\definecolor{rojo}{HTML}{C4142C}
\definecolor{verde}{HTML}{1F9D55}
\definecolor{gris}{HTML}{5D5A73}

\newcolumntype{Y}{>{\raggedright\arraybackslash}X}
\renewcommand{\arraystretch}{1.25}
\setlist{itemsep=1pt,topsep=3pt}
\numberwithin{equation}{section}
\lstset{basicstyle=\ttfamily\small,columns=fullflexible,keepspaces=true,frame=single,rulecolor=\color{gris!40},
  literate={á}{{\'a}}1 {é}{{\'e}}1 {í}{{\'i}}1 {ó}{{\'o}}1 {ú}{{\'u}}1 {ñ}{{\~n}}1 {Ω}{{$\Omega$}}1
    {θ}{{$\theta$}}1 {λ}{{$\lambda$}}1 {Δ}{{$\Delta$}}1 {μ}{{$\mu$}}1 {·}{{$\cdot$}}1 {←}{{$\leftarrow$}}1 {§}{{\S}}1 {→}{{$\to$}}1 {≠}{{$\neq$}}1 {É}{{\'E}}1 {Á}{{\'A}}1 {∨}{{$\vee$}}1 {∧}{{$\wedge$}}1 {σ}{{$\sigma$}}1 {×}{{$\times$}}1 {≈}{{$\approx$}}1}

\newcommand{\rev}[1]{\widetilde{#1}}
\newcommand{\I}{\mathbf{I}}
\newcommand{\vect}[1]{\bm{#1}}
\newcommand{\A}{\mathcal{A}}

\newtcolorbox{meta}{colback=morado,colframe=morado,coltext=white,arc=3mm}
\newtcolorbox{lectura}{colback=morado!5,colframe=morado,fonttitle=\bfseries,title={\faBook~Lecturas}}
\newtcolorbox{ejercicio}[1]{colback=amarillo!10,colframe=amarillo!80!black,fonttitle=\bfseries,title={\faPencil*~Ejercicio #1}}
\newtcolorbox{control}{colback=rojo!5,colframe=rojo,fonttitle=\bfseries,title={\faQuestionCircle~Preguntas de control}}
\newtcolorbox{codigo}[1][De la ecuación al código]{colback=verde!6,colframe=verde,fonttitle=\bfseries,title={\faCode~#1}}
\newtcolorbox{cierre}{colback=verde!6,colframe=verde,fonttitle=\bfseries,title={\faFlagCheckered~Cierre del módulo}}
\newtcolorbox{nota}[1]{colback=amarillo!12,colframe=amarillo!80!black,fonttitle=\bfseries,title={\faLightbulb~#1}}

% \cabecera{número}{título}{subtítulo}
\newcommand{\cabecera}[3]{%
  \title{\textcolor{morado}{\textbf{Ficha #1. #2}}\\[3pt]\large Taller virtual ATA EPE · #3}
  \author{}\date{}}

\cabecera{8}{Guardar, compartir y llevar a la mesa}{Módulo 8: compartir}

\begin{document}
\maketitle
\vspace{-2.5em}

\begin{meta}
\textbf{\textcolor{amarillo}{Meta del módulo.}} Que un montaje se pueda guardar, abrir en otro aparato y reproducir
igual; compartirlo con un enlace; proponer retos (armar los siete prototipos de la cartilla) que el simulador sabe
verificar; y exportar lo armado, con su movimiento, a realidad aumentada. Cara pública: \textbf{compartir y retos}.
\end{meta}

\begin{lectura}
\begin{enumerate}[leftmargin=*]
  \item \textbf{Especificación glTF 2.0} (Khronos Group): nodos, mallas y animaciones; el formato que ya usa la web.
  \item \textbf{Goldberg, «What every computer scientist should know about floating-point arithmetic»},
        \emph{ACM Computing Surveys} 23(1), 1991: por qué el determinismo es delicado.
  \item Documentación de \texttt{<model-viewer>} (realidad aumentada) y de \texttt{CompressionStream} (navegador).
  \item Lo hecho en \texttt{herramientas/generar\_mecanismos.py}: ya muestrea animaciones a 30 cuadros por segundo y
        las exporta a GLB y USDZ.
\end{enumerate}
\end{lectura}

% =================================================================================================
\section{El montaje como documento}

Un montaje se guarda como su \textbf{grafo de ensamble}, no como el estado de la simulación: piezas (catálogo y
motor), uniones (qué tornillo, qué agujeros, estado de la tuerca), bandas y parámetros. Con un número de versión, para
que los montajes viejos se sigan abriendo cuando el formato cambie.

\begin{lstlisting}
{ "version": 1,
  "piezas": [ { "id": "viga-9", "motor": [1, 0, 0, ...] }, ... ],
  "uniones": [ { "tornillo": 3, "piezas": [0, 1], "tuerca": "floja" }, ... ],
  "bandas": [ { "poleas": [4, 5], "cruzada": false } ] }
\end{lstlisting}

\begin{ejercicio}{8.1}
Diseña el formato completo. ¿Qué información se puede \emph{reconstruir} al abrir (masas, cuerpos compuestos, formas
de colisión) y por tanto no se guarda? ¿Cómo se valida un archivo dañado o malicioso?
\end{ejercicio}

% =================================================================================================
\section{Determinismo}

Mismo montaje, mismo paso de tiempo y mismo orden de operaciones dan el mismo resultado \emph{en el mismo motor de
JavaScript}. Entre aparatos hay una trampa: la norma IEEE 754 garantiza que $+$, $-$, $\times$, $\div$ y la raíz
cuadrada den exactamente el mismo resultado en todas partes, pero \emph{no} las funciones trascendentes
(\texttt{Math.sin}, \texttt{Math.exp}), que pueden diferir en el último bit. En un sistema caótico (una biela que se
traba o no), esa diferencia crece.

\begin{ejercicio}{8.2}
Propón dos soluciones: (a) implementar seno, coseno y exponencial propios con solo las operaciones garantizadas;
(b) guardar «fotos» del estado cada cierto tiempo. ¿Qué cuesta cada una? ¿Cuál basta para compartir un montaje?
\end{ejercicio}

% =================================================================================================
\section{Compartir con un enlace}

El documento JSON se comprime (\texttt{CompressionStream('deflate-raw')}), se codifica en base64 apto para URL y va en
el fragmento del enlace (\texttt{\#m=...}), que no se envía al servidor. Un montaje típico ocupa unos cientos de bytes:
cabe en un código QR.

% =================================================================================================
\section{Retos verificables}

Cada reto es un predicado sobre el montaje, calculado con lo construido en los módulos anteriores:
\begin{itemize}[leftmargin=*]
  \item «Arma la biela-manivela»: un grado de libertad (módulo 3) y la trayectoria del pescador sube y baja.
  \item «Haz que la rueda grande gire tres veces más lento»: relación de transmisión $1/3$ medida (módulo 6).
  \item «Equilibra la balanza con pesas distintas»: $m_1x_1 = m_2x_2$ por trabajo virtual (módulo 7).
\end{itemize}
Un reto no compara con una solución guardada: comprueba la física. Así hay muchas formas de resolverlo.

% =================================================================================================
\section{A la mesa: realidad aumentada}

Se simula el montaje unos segundos, se muestrea el motor de cada pieza a 30 cuadros por segundo y se escribe un GLB
con una pista de animación por pieza (traslación y rotación de cada motor), como hace
\texttt{generar\_mecanismos.py}. En Android, \texttt{<model-viewer>} lo muestra en la mesa con movimiento. En iPhone,
Quick Look necesita USDZ, y el exportador del navegador no incluye animación: por ahora el montaje se ve quieto.

\begin{nota}{Bucle sin salto}
Para que la animación se repita sin salto, se muestrea un número entero de periodos del mecanismo: el periodo de la
manivela, o el mínimo común múltiplo cuando hay varias ruedas (como en los mecanismos de la web).
\end{nota}

\begin{codigo}[Qué programar]
\begin{tabularx}{\linewidth}{@{}l Y@{}}
\texttt{app/archivo.js} & Guardar y abrir montajes (JSON con versión y validación); enlace comprimido; QR.\\
\texttt{app/retos.js} & Retos como predicados físicos; los siete prototipos de la cartilla.\\
\texttt{app/exportar.js} & Muestreo de la simulación y GLB animado; visor de realidad aumentada.\\
\texttt{pruebas/} & Guardar y abrir da el mismo grafo; misma simulación en el mismo motor; retos resueltos y no
  resueltos.\\
\end{tabularx}
\end{codigo}

\begin{control}
\begin{enumerate}[leftmargin=*]
  \item ¿Por qué se guarda el grafo de ensamble y no el estado de la simulación?
  \item ¿Qué garantiza IEEE 754 y qué no, y por qué importa en un sistema caótico?
  \item ¿Por qué un reto debe comprobar la física y no comparar con una solución?
  \item ¿Cómo se consigue que la animación exportada se repita sin salto?
\end{enumerate}
\end{control}

\begin{cierre}
Ejercicios 8.1--8.2 · pruebas de archivo, enlace y retos · los siete prototipos de la cartilla armados y verificados
en el taller · cara pública: compartir, retos y realidad aumentada. \textbf{Fin del curso.}
\end{cierre}

\end{document}
